\section{Experiments}
\label{experiments}

\subsection{Research Questions}
\label{sec:exp-rqs}

The evaluation is organized around the following research questions:
\begin{itemize}[leftmargin=*]
  \item RQ1: How accurately does the workflow align natural-language requirements with architectural elements?
  \item RQ2: How complete and useful are the generated requirement-to-architecture trace links?
  \item RQ3: How much do architecture knowledge and agent review contribute to mapping quality?
  \item RQ4: How do human reviewers assess the usefulness of the alignment artifacts?
\end{itemize}

\subsection{Experimental Setup}
\label{sec:exp-setup}

This subsection will describe the implementation setting, agent roles, knowledge sources, prompting protocol, human review procedure, and evaluation process.
It should report configuration details without overstating tool generality.

\subsection{Datasets}
\label{sec:exp-datasets}

The dataset should contain requirement documents paired with reference or expert-reviewed architectural artifacts.
The paper should describe project domains, requirement counts, artifact types, and annotation procedures once the dataset is finalized.

\subsection{Baselines}
\label{sec:exp-baselines}

Baselines should reflect requirement-to-architecture mapping alternatives, such as direct single-agent mapping, LLM-only mapping without architecture knowledge, and mapping without reviewer feedback.

\subsection{Evaluation Metrics}
\label{sec:exp-metrics}

\subsubsection{Architecture Alignment Accuracy}
\label{sec:exp-alignment-accuracy}

Measures whether generated architectural elements and constraints match reference or expert-accepted mappings.

\subsubsection{Requirement Traceability Quality}
\label{sec:exp-traceability-quality}

Measures the correctness, usefulness, and evidence support of requirement-to-architecture trace links.

\subsubsection{Mapping Completeness}
\label{sec:exp-mapping-completeness}

Measures whether architecturally significant requirements are covered by components, responsibilities, constraints, or trace links.

\subsubsection{Architecture Consistency}
\label{sec:exp-architecture-consistency}

Measures whether the produced architecture mapping avoids unsupported decisions, responsibility conflicts, and constraint violations.

\subsection{Main Results}
\label{sec:exp-main-results}

This subsection will report quantitative and qualitative findings after the evaluation is conducted.
No result claims should be added before experimental evidence is available.

\subsection{Ablation Study}
\label{sec:exp-ablation}

The ablation study should isolate the contribution of architecture knowledge, role-separated agents, reviewer checks, and human feedback.

\subsection{Case Study}
\label{sec:exp-case-study}

The case study should illustrate one requirement-to-architecture mapping workflow, including requirement evidence, architectural concerns, mapped components, constraints, and reviewer feedback.
